\chapter{Historie}
\label{chapter:history}

%Als erstes Computerspiel der Geschichte gilt das 1958 von William Higinbotham entwickelte Tennis for Two.

Die ersten Spiele, in denen man auf virtuelle Gegner schie�en konnte, waren sog. Shot'em'Ups, Spiele in 2D-Grafik, bei
denen der Spieler Gegnerwellen, die langsam auf ihn zukommen, abschie�en muss. Besonders Space Invaders aus dem Jahr
1978 erreichte gro�e Beliebtheit. Gewaltdarstellungen waren zu dieser Zeit noch sehr abstrakt. So l�sten sich zum
Beispiel getroffene Alien-Angreifer lediglich mit einer kleinen flimmernden Pixel-Wolke in Luft auf. 

Der erste First-Person-Shooter im heutigen Sinne wurde 1973 von dem Studenten Steve Colley entwickelt. In Maze War
konnte man sich bereits in einer 3D-Welt bewegen und auf Gegner schie�en, sowie �ber ein Netzwerk mit ca. 30 Spielern
zusammen spielen. Dennoch war es so gut wie gewaltfrei: Treffer wurden durch blinken des Bildschirms (oder Teilen davon)
erkenntlich gemacht, getroffene Spieler wurden lediglich an eine zuf�llige Position bef�rdert und von dort aus
weiterspielen.

Ein weiterer Meilenstein in der Entwicklung von Computerspielen und gleichzeitig ein erstes negatives Extrem im Bezug
auf Gewaltdarstellungen war Wolfenstein 3D, das 1992 von id Software ver�ffentlicht wurde. Wolfenstein 3D wurde von der
Bundespr�fstelle f�r jugendgef�hrdende Schriften 1994 aufgrund der sehr realistischen Darstellung der T�tung
menschlicher Gegner, sowie, weil der Spieler gezwungen war, massenhaft menschliche Gegner zu t�ten, um selbst zu
�berleben, indiziert. \\

\begin{figure}[htbp]
\centering
\includegraphics[width=0.8\textwidth]{images/uskfreigaben2012}
\caption{Statistik der USK}
\label{uskfreigaben2012}
\end{figure}

Die Statistik der USK zeigt, dass heute ebenfalls viele Spiele f�r Jugendliche ungeeignete Darstellungen von
Gewalt enthalten. 2012 erhielten 8,7 Prozent der gepr�ften Spiele keine Freigabe f�r Jugendliche.
